% Generated corrected-evidence block; see registry and audit.
\subsection{Day-in-the-Life Operations}\label{sec:sim-leop}
The corrected exports under \path{evidence/new-science-ditl/} replace the previously ingested incorrect/incomplete Science/DITL evidence. The old raw files, analyses and PDFs remain archived and must not support current numerical conclusions. The registry selects one source identity per scenario; no rows are pooled across runs. All three exports reach their named timestamp horizons: two 24-hour cases and one 365-day case. Absolute contact and SAA timestamps consistently identify 1 March 2000 UTC, not the approved 2029 analysis epoch.

The files contain no producing commit or runtime manifest. RAAN, mass, battery capacity, integrator, selected hardware and elevation mask therefore remain unverified. The two day cases have matching recorded orbital positions and contact windows to numerical precision, but complete configuration equivalence is not established. Their comparison is a sensitivity study, not a controlled causal test of LEOP alone.

The following telemetry integrals hold each recorded state over $[t_i,t_{i+1})$, exclude terminal-row duration, and sum to the recorded horizon. These exports lack explicit interval-end fields. Transition-log mode durations and slew-event durations are reported separately in Appendix~\ref{app:ditl-detail}; disagreements are not silently removed. Battery percentages elsewhere use 25 Wh only as an identified planning normalization, not a recovered capacity.
\begin{table}[H]\centering\small
\caption{Corrected telemetry accounting; left-held interval interpretation, historical epoch}\label{tab:annual-ditl}
\begin{tabularx}{\linewidth}{@{}Xrrr@{}}\toprule
Metric & LEOP 24 h & No-LEOP 24 h & LEOP 365 d\\\midrule
Recorded timestamp span (s) & 86,400 & 86,400 & 31,536,000\\
Sun Pointing (s) & 79,268 & 81,852 & 29,660,918\\
Detection (s) & 2,676 & 2,674 & 894,054\\
Relay (s) & 894 & 894 & 269,876\\
Slewing (s) & 980 & 980 & 375,688\\
Safe (s) & 0 & 0 & 332,882\\
Geographic SAA (s) & 2,696 & 2,696 & 881,882\\
Acquisition gate (s) & 758 & 758 & 247,258\\
Minimum recorded battery (Wh) & 15.128 & 18.004 & 15.098\\
Recorded solar integral (Wh) & 145.575 & 150.403 & 50,288.809\\
Recorded load integral (Wh) & 85.167 & 85.719 & 30,866.103\\
Native generation (source MB) & 0.009 & 0.009 & 3.514\\
Native transfer (source MB) & 0.009 & 0.009 & 3.514\\
\bottomrule\end{tabularx}\end{table}
\begin{figure}[H]\centering\includegraphics[width=\linewidth]{evidence/science-ditl/figures/corrected_20261109_day_occupancy_battery.png}
\caption{Corrected day-run mode occupancy and battery normalized to the 25 Wh planning reference. Same time span does not establish identical runtime configuration.}\label{fig:annual-mode-exposure}\end{figure}
Both day runs record 2,696 s geographic SAA dwell and 758 s acquisition gate. Removing startup does not increase this gate total in the supplied pair. LEOP includes one 2,460 s detumble event; this is not Monte Carlo reliability. For the annual trace, acquisition/SAA is 28.038\%, with numerator 247,258 s and denominator 881,882 s. Neither fraction is calibrated current-FPDU science duty.

Geographic SAA passages, sampled cell visits, gated exposure, model expected counts, accepted detector events and received science products are distinct. The legacy response does not qualify the six-scintillator FPDU. No MMSC/MSC/EMSC success percentage follows from these operational records.
\subsection{Ground Station and Data Return}\label{sec:sim-coverage}
The corrected annual contact log contains 820 intervals totaling 244,016 s. Mean daily contact is 668.537 s and minimum daily contact 514 s; the longest intercontact gap is 48,458 s. The source mask is unknown and is not applied a second time. Visibility alone corresponds to 585,638,400 bytes at 19.2 kbps. Intersecting these contacts with telemetry-labelled Relay yields 178,682 s, or 428,836,800 ideal bytes. Transition-log Relay instead gives 428,836,800 bytes; the timing-definition sensitivity remains explicit.

Native annual transfer credit is 3.514083862 source-labelled MB, separate from hypothetical v3 delivery. Source MB scaling and received-packet integrity are not demonstrated by the label. The v3 preferred EMSC workload generates 46,637,364 bytes with peak modeled queue 75,336 bytes and final backlog 0 bytes. It uses corrected, time-aligned sampled gates and Relay visibility at the full ideal 2400 bytes/s; it is not emitted v3 telemetry. Appendix~\ref{app:science-data-budget} gives tier and outage sensitivity.
The preliminary communications assessment uses the full idealized 19.2 kbps gross transfer opportunity during eligible ground-station relay windows. No arbitrary link derating is imposed at MCR. Detailed RF link-budget closure, protocol overhead, packet losses, retransmissions, and achievable net payload throughput shall be assessed at PDR.
